\section{Services, foundation, platform}
\label{sec:services}

\subsection{Services}
Cross-cutting facilities used by Engine/Game/Rendering:
\begin{itemize}
  \item \textbf{Logging} --- \symbolref{Logger} (file + terminal; level from env).
  \item \textbf{Config} --- \symbolref{EnvVars} / \symbolref{IEnvVars}
        (defaults filled in \symbolref{Illumo::Init}: canvas size, fps/tps,
        synchronized presentation via \texttt{vsync}, mode string, \ldots).
        The production file is \texttt{envvars.json}
        beside the executable, seeded from the tracked default only when absent;
        it never depends on the launch working directory.
  \item \textbf{Input} --- \symbolref{InputManager}, \symbolref{InputContext},
        \symbolref{KeyCode}; action binding used by the game module.
        No Game types (D-E2).
  \item \textbf{CLI} --- \symbolref{CommandLine} owns the token UI and batched
        geometry. \symbolref{CommandRegistry} owns command metadata and
        callbacks; \symbolref{SysCmdLine} handles startup options. Help,
        version, and window/canvas dimensions are applied before window creation.
        General tooling stays in Services; \symbolref{CellGameModule} registers
        simulation/canvas/camera/save-load commands with usage, descriptions,
        and completion candidates. Editing includes selection, word movement and
        deletion, history, quoted arguments, measured caret placement, command
        chaining with \texttt{;}, alias macro management (\texttt{alias}/\texttt{unalias}),
        inline ghost-text auto-suggestions, dynamic parameter hints, and a
        horizontally scrolling input viewport. The current history draw path
        truncates each entry to panel width and scrolls by raw history entries;
        D-UI2's word-wrap and visual-line metrics are not currently
        implemented. The easy-font batch is sized for 8,000 UI quads. The
        console is Debug-only product composition
        (D-B1, D-CLI1, D-UI1--D-UI3).
  \item \textbf{Allocators} --- pool/arena/debug allocators (experimental /
        learning surface; not all wired into the hot path).
  \item \textbf{SaveLoad API} --- pure declarations; OS dialogs implemented
        per platform.
\end{itemize}

\subsection{Foundation}
Headers that should stay cheap to include:
macros (\pathref{MacroDefs.h}), build info, \pathref{MathTypes.h},
\pathref{ArrayQueue.h}.

\begin{statusbox}
\pathref{MacroDefs.h} still pulls Windows headers on \texttt{\_WIN32} (for
\texttt{Sleep} / \texttt{\_\_SLEEP}). That is include-heavy for a widely used
header. \textbf{Deferred (D-F1):} leave as-is until compile time, name
collisions, or header hygiene actually hurt; then split platform sleep into a
narrower header.
\end{statusbox}

\subsection{Platform}
Thin platform sources:
\begin{itemize}
  \item Windows: \pathref{WinMain.cpp}, \pathref{WinSaveLoad.cpp} ---
        supported and currently verified.
  \item Linux: \pathref{\_main.cpp}, \pathref{LinuxSaveLoad.cpp}
        --- unsupported stale scaffold; the entry references removed shared
        bootstrap APIs.
  \item macOS: \pathref{Main.cpp}, \pathref{MacSaveLoad.mm}
        --- unsupported stale scaffold; the entry references removed APIs and
        CMake does not enable Objective-C++ for the selected \texttt{.mm} file.
\end{itemize}

Contract: bootstrap $\rightarrow$ call \symbolref{CellMain}. Each port also
implements the native load-location and save-location functions in
\symbolref{SaveLoad}.
Source presence and CMake selection do not establish support; Linux/macOS
require native configure, build, tests, rendering/input, dialogs, and shutdown
validation.

\subsection{Third-party}
Vendored under \pathref{Illumo/thirdparty/}. Policy from \pathref{docs/contributing.md}:
new third-party deps require review/PR. Existing: GLFW, GLEW, GLM, FreeType,
JSON, STB, Tracy, image libs, \ldots
